% Options for packages loaded elsewhere
\PassOptionsToPackage{unicode}{hyperref}
\PassOptionsToPackage{hyphens}{url}
%
\documentclass[
]{article}
\usepackage{amsmath,amssymb}
\usepackage{iftex}
\ifPDFTeX
  \usepackage[T1]{fontenc}
  \usepackage[utf8]{inputenc}
  \usepackage{textcomp} % provide euro and other symbols
\else % if luatex or xetex
  \usepackage{unicode-math} % this also loads fontspec
  \defaultfontfeatures{Scale=MatchLowercase}
  \defaultfontfeatures[\rmfamily]{Ligatures=TeX,Scale=1}
\fi
\usepackage{lmodern}
\ifPDFTeX\else
  % xetex/luatex font selection
\fi
% Use upquote if available, for straight quotes in verbatim environments
\IfFileExists{upquote.sty}{\usepackage{upquote}}{}
\IfFileExists{microtype.sty}{% use microtype if available
  \usepackage[]{microtype}
  \UseMicrotypeSet[protrusion]{basicmath} % disable protrusion for tt fonts
}{}
\makeatletter
\@ifundefined{KOMAClassName}{% if non-KOMA class
  \IfFileExists{parskip.sty}{%
    \usepackage{parskip}
  }{% else
    \setlength{\parindent}{0pt}
    \setlength{\parskip}{6pt plus 2pt minus 1pt}}
}{% if KOMA class
  \KOMAoptions{parskip=half}}
\makeatother
\usepackage{xcolor}
\usepackage[margin=1in]{geometry}
\usepackage{longtable,booktabs,array}
\usepackage{calc} % for calculating minipage widths
% Correct order of tables after \paragraph or \subparagraph
\usepackage{etoolbox}
\makeatletter
\patchcmd\longtable{\par}{\if@noskipsec\mbox{}\fi\par}{}{}
\makeatother
% Allow footnotes in longtable head/foot
\IfFileExists{footnotehyper.sty}{\usepackage{footnotehyper}}{\usepackage{footnote}}
\makesavenoteenv{longtable}
\usepackage{graphicx}
\makeatletter
\newsavebox\pandoc@box
\newcommand*\pandocbounded[1]{% scales image to fit in text height/width
  \sbox\pandoc@box{#1}%
  \Gscale@div\@tempa{\textheight}{\dimexpr\ht\pandoc@box+\dp\pandoc@box\relax}%
  \Gscale@div\@tempb{\linewidth}{\wd\pandoc@box}%
  \ifdim\@tempb\p@<\@tempa\p@\let\@tempa\@tempb\fi% select the smaller of both
  \ifdim\@tempa\p@<\p@\scalebox{\@tempa}{\usebox\pandoc@box}%
  \else\usebox{\pandoc@box}%
  \fi%
}
% Set default figure placement to htbp
\def\fps@figure{htbp}
\makeatother
\setlength{\emergencystretch}{3em} % prevent overfull lines
\providecommand{\tightlist}{%
  \setlength{\itemsep}{0pt}\setlength{\parskip}{0pt}}
\setcounter{secnumdepth}{5}
\usepackage{float}
\renewcommand{\thefigure}{S\arabic{figure}}
\renewcommand{\thetable}{S\arabic{table}}
\usepackage{booktabs}
\usepackage{longtable}
\usepackage{array}
\usepackage{multirow}
\usepackage{wrapfig}
\usepackage{float}
\usepackage{colortbl}
\usepackage{pdflscape}
\usepackage{tabu}
\usepackage{threeparttable}
\usepackage{threeparttablex}
\usepackage[normalem]{ulem}
\usepackage{makecell}
\usepackage{xcolor}
\usepackage{bookmark}
\IfFileExists{xurl.sty}{\usepackage{xurl}}{} % add URL line breaks if available
\urlstyle{same}
\hypersetup{
  pdftitle={Supplementary Results},
  pdfauthor={Amber M. Young},
  hidelinks,
  pdfcreator={LaTeX via pandoc}}

\title{Supplementary Results}
\usepackage{etoolbox}
\makeatletter
\providecommand{\subtitle}[1]{% add subtitle to \maketitle
  \apptocmd{\@title}{\par {\large #1 \par}}{}{}
}
\makeatother
\subtitle{Survival-guided matrix factorization identifies reproducible
prognostic programs in pancreatic cancer}
\author{Amber M. Young}
\date{}

\begin{document}
\maketitle

{
\setcounter{tocdepth}{2}
\tableofcontents
}
\newpage
\phantomsection
\addcontentsline{toc}{section}{List of Figures}
\listoffigures
\newpage
\phantomsection
\addcontentsline{toc}{section}{List of Tables}
\listoftables
\newpage

\section{The DeSurv model shows consistent convergence and stability
across
restarts}\label{the-desurv-model-shows-consistent-convergence-and-stability-across-restarts}

Across datasets and initialization schemes, the DeSurv algorithm
consistently converged to numerically stable solutions within the
designated iteration budget. Warm-start strategies, in which solutions
at \(\alpha=0\) initialized supervised runs, substantially reduced
variability across restarts and improved reproducibility. Fig.
\ref{fig:fig-converge} shows representative loss trajectories
demonstrating monotone decreases until convergence across restarts.
Compared with naive random initialization, warm-starts produced tighter
distributions of cross-validated C-index and partial likelihood,
confirming stability of the optimization procedure.

\begin{figure*}[htbp]
\includegraphics[width=5in,height=3in]{figures/fig-converge-1} \caption{Convergence of model across initializations. \label{fig:fig-converge}}\label{fig:fig-converge}
\end{figure*}

\section{Simulation results under null and mixed
scenarios}\label{simulation-results-under-null-and-mixed-scenarios}

To characterize the operating characteristics of DeSurv across different
signal regimes, we evaluated performance under three simulation
scenarios. In the primary scenario (main text Fig. 3), prognostic
programs explained low variance relative to outcome-neutral background
signals, and survival depended on 150 factor-specific marker genes
(\(\beta_1 = 2\), \(\beta_2 = \beta_3 = 0\)). In the null scenario,
survival times were generated independently of gene expression
(\(\beta = 0\)), so no prognostic structure existed. In the mixed
scenario, survival depended on 300 genes per lethal factor, of which
50\% were factor-specific markers and 50\% were background genes with
loadings shared across all factors. All scenarios used \(G = 3{,}000\)
genes, \(N = 200\) samples, \(K = 3\) true factors, and identical gamma
distribution parameters for the gene loading matrix \(W\) (see Materials
and Methods). Performance was evaluated using cross-validated
concordance index (C-index), precision of recovered prognostic gene
sets, and the distribution of selected ranks across simulation
replicates.

Under the null scenario (Fig. \ref{fig:sim-null-mixed}A--B), DeSurv
showed no advantage over unsupervised NMF (\(\alpha = 0\)). Both methods
yielded C-index values near 0.5, consistent with the absence of survival
signal. The distribution of selected ranks was concentrated at low
values for both methods, with NMF more strongly favoring \(k = 2\).
These results confirm that DeSurv does not hallucinate prognostic
structure when none exists: Bayesian optimization selected low values of
the supervision parameter \(\alpha\), effectively recovering standard
NMF as a special case. This specificity control is important because it
establishes that the advantages observed in the primary scenario reflect
genuine signal recovery rather than overfitting to noise in the survival
labels.

Under the mixed scenario (Fig. \ref{fig:sim-null-mixed}C--F), four
mechanistic observations characterize DeSurv's advantage. First, when
precision is evaluated against the factor-specific marker genes alone
(150 genes; Fig. \ref{fig:sim-null-mixed}D left), DeSurv substantially
outperforms NMF---matching or exceeding primary-scenario levels---even
though overall survival-gene precision (300-gene survival set including
background genes; Fig. \ref{fig:sim-null-mixed}D right) is more modestly
improved. This dissociation reveals that background-gene recovery is
similar between methods, and DeSurv's edge comes specifically from
concentrating marker programs. Second, the learned factor most aligned
with the true marker program receives a substantially larger Cox
coefficient (\(|\beta|\)) under DeSurv than under NMF (Fig.
\ref{fig:sim-null-mixed}E), confirming that survival supervision
preferentially amplifies the marker-specific signal. Third, C-index
values are modestly higher for DeSurv (Fig. \ref{fig:sim-null-mixed}C),
consistent with the moderate overall improvement when background-gene
dilution attenuates the survival signal. Together, these results
demonstrate that DeSurv's advantage in the mixed scenario is
mechanistic: the method specifically recovers marker programs and
assigns them elevated survival weights, rather than incidental capture
of background-gene prognosis.

Together, the three scenarios provide a dose-response relationship:
DeSurv's advantage is greatest when variance and prognosis diverge
(primary), moderate when they partially overlap (mixed), and absent when
no survival signal exists (null). The mechanistic basis for the
attenuation in the mixed scenario is now clear: DeSurv continues to
specifically recover marker programs (high marker precision, elevated
\(|\beta|\)), but the overall survival-gene precision improvement is
dampened because background-gene dilution in the 300-gene survival
signal reduces the discrimination advantage. This gradient is consistent
with the predictions of sufficient dimension reduction theory and offers
practical guidance for users: DeSurv is most beneficial in settings
where the dominant transcriptional axes are expected to be
prognostically neutral, as is common in cancers with high tumor purity
variation or dominant tissue-composition signals.

\begin{figure*}[t]
\includegraphics[width=\textwidth,height=6.5in]{figures/fig-sim-null-mixed-1} \caption{DeSurv performance under null and mixed simulation scenarios. (A--B) Null scenario ($\beta = 0$): both methods yield C-index $\approx 0.5$ (A) and comparable rank distributions (B), confirming DeSurv does not hallucinate prognostic structure; precision is omitted as no true survival genes exist. (C--F) Mixed scenario: survival depends on factor-specific marker genes (50\%) and shared background genes (50\%). (C) C-index is modestly higher for DeSurv. (D) Precision by gene-set definition: marker-only precision (150 genes; left facet) substantially favors DeSurv; full 300-gene survival-set precision (right facet) shows a more modest improvement. (E) The matched factor receives a substantially larger Cox coefficient $|\beta|$ under DeSurv, confirming preferential amplification of marker-specific signal. (F) Selected-rank distributions. DeSurv (blue); NMF (red, $\alpha = 0$); both Bayesian-optimized. \label{fig:sim-null-mixed}}\label{fig:fig-sim-null-mixed}
\end{figure*}

\section{Cross-validated C-index across factorization
rank}\label{cross-validated-c-index-across-factorization-rank}

To evaluate how the number of latent factors \(k\) affects prognostic
performance, we performed an exhaustive grid search over
\(k \in \{2, 3, \ldots, 12\}\) and supervision strength
\(\alpha \in \{0, 0.05, \ldots, 0.95\}\) using 5-fold cross-validation
on the TCGA+CPTAC training cohort. For each \(k\), the best-performing
\(\alpha\) was selected by maximizing the mean CV C-index. The
unsupervised NMF baseline (\(\alpha = 0\)) was evaluated at the same
\(k\) values for comparison. Results are shown for two gene-selection
settings: \(n_{\text{top}} = 270\) top genes per factor and
\(n_{\text{top}} = \text{ALL}\) (all genes retained), both at
\(\lambda = 0.349\) and \(\xi = 0.056\).

Fig. \ref{fig:fig-cindex-by-k} shows the resulting C-index curves for
both training (CV) and external validation (pooled across held-out
cohorts), with shaded ribbons indicating \(\pm 1\) standard error. In
training, DeSurv consistently outperforms NMF across all \(k\),
confirming that survival supervision improves in-sample concordance. In
external validation, DeSurv at \(k = 3\) achieves concordance comparable
to NMF at higher ranks (e.g., \(k = 7\); Table
\ref{tab:tab-cindex-comparison}), demonstrating that survival
supervision concentrates prognostic information into fewer factors
rather than requiring additional rank to capture it.

\begin{figure*}[t]
\includegraphics[width=\textwidth]{figures/fig-cindex-by-k-1} \caption{Training (CV) and external validation C-index as a function of factorization rank $k$ for DeSurv (blue) and standard NMF (red, $\alpha = 0$). Shaded ribbons show $\pm 1$ SE. For each $k$, the best $\alpha$ was selected by maximizing the mean CV C-index. Validation C-index is pooled across all external cohorts. (A) Results using $n_{\text{top}} = 270$ top genes per factor ($\lambda = 0.349$, $\xi = 0.056$). (B) Results retaining all genes ($n_{\text{top}} = \text{ALL}$; $\lambda = 0.349$, $\xi = 0.056$).}\label{fig:fig-cindex-by-k}
\end{figure*}

\section{Standard NMF rank selection
heuristics}\label{standard-nmf-rank-selection-heuristics}

Standard unsupervised rank selection criteria --- reconstruction
residuals, cophenetic correlation coefficient, and mean silhouette width
--- yielded inconsistent guidance for selecting the factorization rank
\(k\) in the PDAC training cohort (TCGA + CPTAC). Reconstruction
residuals decreased smoothly without a clear elbow (Fig.
\ref{fig:nmf-diagnostics}A). Cophenetic correlation fluctuated without a
distinct transition point (Fig. \ref{fig:nmf-diagnostics}B). Mean
silhouette width favored very low ranks that conflicted with other
criteria (Fig. \ref{fig:nmf-diagnostics}C). These contradictions
motivated the survival-supervised rank selection approach described in
the main text.

\begin{figure*}[t]
\includegraphics[width=\textwidth,height=3.5in]{figures/fig-nmf-diagnostics-1} \caption{Standard unsupervised NMF rank selection heuristics for the PDAC training cohort (TCGA + CPTAC). (A) Reconstruction residuals as a function of factorization rank $k$. (B) Cophenetic correlation coefficient as a function of $k$. (C) Mean silhouette width across multiple distance metrics as a function of $k$. These criteria yield inconsistent guidance, motivating the survival-supervised model selection approach (main text Fig. 2D). \label{fig:nmf-diagnostics}}\label{fig:fig-nmf-diagnostics}
\end{figure*}

\section{K-sensitivity analysis: effect of gene focusing on iCAF factor
validation}\label{k-sensitivity-analysis-effect-of-gene-focusing-on-icaf-factor-validation}

To evaluate the robustness of the K=3 iCAF factorization, we conducted
an exhaustive sensitivity analysis across \(K \in \{2, 3, 5, 7, 9\}\)
and \(\alpha \in \{0, 0.25, 0.35, 0.55, 0.75, 0.85, 0.95\}\) (35
combinations), holding all remaining hyperparameters fixed at the
production-optimal values (\(\lambda = 0.349\), \(\xi = 0.056\), 100
initializations). For each fit, we identified the factor most correlated
with the production K=3 D1 (the iCAF factor) and evaluated its
prognostic value in 570 independent validation patients across 4 cohorts
(Dijk, Moffitt GEO, Puleo, PACA-AU), using stratified Cox regression
(\texttt{strata(dataset)}) with and without adjustment for PurIST and
DeCAF classifiers. Factor-program fidelity was quantified by Spearman
rank correlation of \(n_{\text{top}} = 270\)-focused H-scores (H-cor)
with the production factor. We repeated this analysis under two
gene-projection settings, \(n_{\text{top}} = 270\) (production-optimal
top genes per factor) and \(n_{\text{top}} = \text{NULL}\) (all 1,970
genes), to isolate the effect of gene focusing on validation
performance.

The central finding is shown in Table \ref{tab:tab-ksens-summary}:
\textbf{\(n_{\text{top}} = 270\) is the critical hyperparameter enabling
the production fit's adjusted validation.} Moving from
\(n_{\text{top}} = \text{NULL}\) to \(n_{\text{top}} = 270\), while
holding \(\lambda\) and \(\xi\) at their optimal values, increases the
number of K \(\times\) \(\alpha\) combinations achieving adjusted
\(P < 0.05\) from 5/35 to 18/35 and restores K=3 at \(\alpha=0.55\) to
adjusted \(P = 0.003\), nearly identical to the production fit
(\(P = 0.004\)). By contrast, changing only \(\lambda\) from 0.3 to
0.349 while keeping \(n_{\text{top}} = \text{NULL}\) leaves the overall
count unchanged (5/35), confirming that \(\lambda/\xi\) tuning alone
cannot reproduce the adjusted signal.

\begin{table}[!h]
\centering
\caption{\label{tab:tab-ksens-summary}Three-way comparison of K-sensitivity analyses differing only in gene-projection setting ($n_{\text{top}}$) and regularization ($\lambda$, $\xi$). All analyses span K $\in$ \{2, 3, 5, 7, 9\} $\times$ $\alpha \in$ \{0, 0.25, 0.35, 0.55, 0.75, 0.85, 0.95\} = 35 combinations. Count shows combinations achieving adjusted $P < 0.05$; bold highlights the $n_{\text{top}}=270$ result that restores K=3 adjusted significance.}
\centering
\resizebox{\ifdim\width>\linewidth\linewidth\else\width\fi}{!}{
\fontsize{9}{11}\selectfont
\begin{tabular}[t]{llll}
\toprule
Analysis condition & Adj $P < 0.05$ (of 35) & K=3, $\alpha$=0.55 adj $P$ & Best single combination\\
\midrule
$n_{\text{top}}$ = NULL, $\lambda$ = 0.300, $\xi$ = 0.050 & 5 / 35 & 0.072 & K=7, $\alpha$=0.75 (adj $P$ = 0.019)\\
$n_{\text{top}}$ = NULL, $\lambda$ = 0.349, $\xi$ = 0.056 & 5 / 35 & 0.097 & K=5, $\alpha$=0.25 (adj $P$ = 0.0003)\\
$n_{\text{top}}$ = 270,\hspace{2pt} $\lambda$ = 0.349, $\xi$ = 0.056 & \textbf{18 / 35} & \textbf{0.003} & K=3, $\alpha$=0.55 (adj $P$ = 0.003)\\
\bottomrule
\end{tabular}}
\end{table}

Tables \ref{tab:tab-ksens-adjp-ntop270} and
\ref{tab:tab-ksens-adjp-allgenes} present the full K \(\times\)
\(\alpha\) adjusted P-value matrices for the \(n_{\text{top}} = 270\)
and \(n_{\text{top}} = \text{NULL}\) analyses respectively. With
\(n_{\text{top}} = 270\), K=3 achieves adjusted significance at 5 of 7
\(\alpha\) values tested
(\(\alpha \in \{0.00, 0.25, 0.55, 0.75, 0.95\}\)), and broad adjusted
significance is also observed at K=2 (5/7) and K=5 (4/7). With
\(n_{\text{top}} = \text{NULL}\) (Table
\ref{tab:tab-ksens-adjp-allgenes}), the landscape is sparse and
structurally different: K=3 at \(\alpha=0.55\) narrowly misses
significance (\(P = 0.097\)), and only isolated combinations scattered
across K=3, K=5, K=7, and K=9 achieve adjusted \(P < 0.05\).

\begin{table}[!h]
\centering
\caption{\label{tab:tab-ksens-adjp-ntop270}Adjusted validation P-values (Cox + PurIST + DeCAF + strata(dataset)) across K $\times$ $\alpha$ with $n_{\text{top}}=270$, $\lambda=0.349$, $\xi=0.056$. Each cell gives the P-value for the factor most correlated with the production K=3 iCAF factor. Bold: $P < 0.05$; 18 of 35 combinations are significant.}
\centering
\fontsize{9}{11}\selectfont
\begin{tabular}[t]{llllllll}
\toprule
\multicolumn{1}{c}{ } & \multicolumn{7}{c}{$\alpha$} \\
\cmidrule(l{3pt}r{3pt}){2-8}
  & 0.00 & 0.25 & 0.35 & 0.55 & 0.75 & 0.85 & 0.95\\
\midrule
K=2 & 0.564 & \textbf{0.038} & \textbf{0.038} & \textbf{0.019} & 0.060 & \textbf{0.011} & \textbf{0.002}\\
K=3 & \textbf{0.030} & \textbf{0.030} & 0.519 & \textbf{0.003} & \textbf{0.028} & 0.263 & \textbf{0.009}\\
K=5 & \textbf{5e-4} & \textbf{0.005} & \textbf{8e-4} & \textbf{0.019} & 0.073 & 0.604 & 0.447\\
K=7 & 0.513 & 0.471 & \textbf{6e-4} & \textbf{0.041} & 0.676 & 0.131 & 0.130\\
K=9 & \textbf{3e-4} & 0.299 & 0.243 & \textbf{0.005} & 0.943 & 0.867 & 0.263\\
\bottomrule
\end{tabular}
\end{table}

\begin{table}[!h]
\centering
\caption{\label{tab:tab-ksens-adjp-allgenes}Adjusted validation P-values across K $\times$ $\alpha$ with $n_{\text{top}}$=NULL (all 1970 genes), $\lambda=0.349$, $\xi=0.056$. Layout identical to Table \ref{tab:tab-ksens-adjp-ntop270}. Only 5 of 35 combinations are significant (bold), and K=3 at $\alpha=0.55$ (the production neighborhood) reaches only $P = 0.097$.}
\centering
\fontsize{9}{11}\selectfont
\begin{tabular}[t]{llllllll}
\toprule
\multicolumn{1}{c}{ } & \multicolumn{7}{c}{$\alpha$} \\
\cmidrule(l{3pt}r{3pt}){2-8}
  & 0.00 & 0.25 & 0.35 & 0.55 & 0.75 & 0.85 & 0.95\\
\midrule
K=2 & 0.414 & 0.192 & 0.142 & 0.197 & 0.090 & 0.194 & 0.158\\
K=3 & 0.205 & \textbf{0.024} & 0.280 & 0.097 & 0.112 & 0.076 & 0.111\\
K=5 & 0.461 & \textbf{3e-4} & 0.150 & 0.233 & 0.067 & 0.124 & 0.197\\
K=7 & 0.394 & 0.336 & 0.304 & 0.097 & 0.094 & 0.062 & \textbf{0.021}\\
K=9 & \textbf{0.023} & 0.093 & 0.474 & \textbf{0.022} & 0.106 & 0.127 & 0.122\\
\bottomrule
\end{tabular}
\end{table}

Four additional findings emerge from these analyses. First, K=7's
apparent advantage observed in a parallel analysis using conventional
hyperparameters (\(\lambda=0.3\), \(n_{\text{top}} = \text{NULL}\)),
where K=7 at \(\alpha=0.75\) achieved adjusted \(P = 0.019\), completely
disappears with \(n_{\text{top}} = 270\) (adj \(P = 0.676\), Table
\ref{tab:tab-ksens-adjp-ntop270}). Focused 270-gene projection fragments
the iCAF signal across K=7's factors, indicating that the earlier K=7
advantage depended on the all-genes projection rather than reflecting an
intrinsic property of the 7-factor solution. Second, K=3 achieves
adjusted significance across a broad \(\alpha\) range with
\(n_{\text{top}} = 270\) (5/7 values), demonstrating that the production
result does not depend on the BO-selected \(\alpha = 0.334\)
specifically. Third, K=5 at \(\alpha=0.55\) achieves the highest
factor-program fidelity in the \(n_{\text{top}} = 270\) analysis (H-cor
= 0.903 vs.~K=3's 0.855) and also achieves adjusted significance
(\(P = 0.019\)), providing convergent evidence for the iCAF biology from
a higher-K model; an independent K=5 all-genes fit at \(\alpha=0.25\)
likewise achieves adjusted significance (\(P = 0.0003\), H-cor = 0.91).
Fourth, without gene focusing, high supervision weights
(\(\alpha \geq 0.75\)) produce factors with extreme training
discrimination (HR \(\approx 0.07\)--\(0.1\), LRT \(P \sim 1e-60\)) that
fail external validation, indicating overfitting of the training
survival signal; gene focusing (\(n_{\text{top}} = 270\)) suppresses
this by constraining each factor to its 270 most biologically
characteristic genes.

\section{K=3 factors are partially but asymmetrically nested within
K=7}\label{k3-factors-are-partially-but-asymmetrically-nested-within-k7}

To assess whether the K=3 solution is simply a pruned version of the K=7
solution, we computed pairwise Spearman correlations between the
W-matrix gene loadings of the production K=3 model and a K=7 DeSurv fit
at \(\alpha = 0.35\) (the nearest grid value to the BO-selected optimum
at \(\alpha = 0.334\)), both using \(n_{\text{top}} = 270\),
\(\lambda = 0.349\), \(\xi = 0.056\).

The correspondence was partial but asymmetric (Fig.
\ref{fig:k3-k7-nesting}). Two of three K=3 factors had clear one-to-one
K=7 counterparts: K3\_F1 mapped strongly to K7\_F1 (\(r = 0.77\)) and
K3\_F2 mapped very strongly to K7\_F7 (\(r = 0.88\)). In contrast,
K3\_F3 showed no single dominant K=7 match; instead, it correlated
moderately with four K=7 factors simultaneously (K7\_F2 through K7\_F6;
\(r = 0.57\)--\(0.71\)). K7\_F4 showed weak-to-moderate correlations
with all three K=3 factors (\(r = 0.45\)--\(0.52\)), suggesting it is a
mixed composite with no clear biological counterpart in the K=3
solution.

This fragmentation pattern indicates that K=7 does not simply subdivide
K=3 structure. Rather, it disperses one coherent K=3 prognostic program
across multiple lower-signal components, consistent with the observed
failure of K=7 to achieve adjusted significance in external cohorts
(adjusted \(P = 0.676\); see main text).

\begin{figure*}[t]
\includegraphics[width=\textwidth]{figures/k3-k7-nesting-1} \caption{Factor correspondence between the production K=3 model and a K=7 DeSurv fit ($\alpha = 0.35$, $n_{\text{top}} = 270$), measured by Spearman rank correlation of W-matrix gene loadings across all genes. Two K=3 factors have clear K=7 counterparts (K3\_F1 $\to$ K7\_F1, $r = 0.77$; K3\_F2 $\to$ K7\_F7, $r = 0.88$), while K3\_F3 is distributed broadly across four K=7 factors ($r = 0.57$--$0.71$), indicating fragmentation rather than nesting.}\label{fig:k3-k7-nesting}
\end{figure*}

\section{Cutpoint selection and validation of dichotomized risk
groups}\label{cutpoint-selection-and-validation-of-dichotomized-risk-groups}

To convert the continuous DeSurv linear predictor into a clinically
interpretable binary risk stratification, we selected an optimal z-score
cutpoint via cross-validation on the TCGA+CPTAC training cohort. For
each candidate cutpoint in \(\{-2.0, -1.8, \ldots, 2.0\}\), we computed
the mean absolute log-rank z-statistic across held-out folds and
selected the cutpoint maximizing this metric (Fig.
\ref{fig:fig-cutpoint-km}A). We then applied this cutpoint to external
validation cohorts by standardizing each patient's linear predictor
using the training mean and standard deviation. Fig.
\ref{fig:fig-cutpoint-km}B shows the resulting Kaplan--Meier curves for
the pooled external validation analysis, with a stratified log-rank test
accounting for dataset of origin. Panels C--F display per-cohort
validation KM curves for Dijk, Moffitt, PACA-AU, and Puleo,
respectively. Across all cohorts, the DeSurv-derived high-risk group
showed consistently shorter survival, confirming that the
training-derived cutpoint generalizes to independent datasets.

\begin{figure*}[t]
\includegraphics[width=\textwidth,height=8in]{figures/fig-cutpoint-km-1} \caption{Cutpoint selection and external validation of dichotomized risk groups. (A) Cross-validated mean absolute log-rank z-statistic as a function of z-score cutpoint; red dashed line indicates the selected optimum. (B) Pooled validation Kaplan--Meier curves (stratified log-rank test). (C--D) Per-cohort validation KM curves: (C) Dijk, (D) Moffitt. (E--F) Per-cohort validation KM curves: (E) PACA-AU, (F) Puleo.}\label{fig:fig-cutpoint-km}
\end{figure*}

\section{Overlap of DeSurv risk groups with known molecular
subtypes}\label{overlap-of-desurv-risk-groups-with-known-molecular-subtypes}

To assess whether the DeSurv-derived risk stratification captures
biologically meaningful transcriptional programs, we examined the
overlap between dichotomized risk groups (High vs Low) and two
established PDAC molecular classifiers: PurIST (Basal-like vs Classical)
and DeCAF (proCAF vs restCAF). Fig. \ref{fig:fig-subtype-overlap} shows
the composition of each risk group across all pooled external validation
cohorts, with Fisher's exact test P-values quantifying the association.
The high-risk group was significantly enriched for Basal-like (PurIST)
and proCAF (DeCAF) subtypes, consistent with the known poor prognosis of
these molecular classes. This concordance confirms that DeSurv's
survival-driven factorization recovers clinically relevant biology
without requiring subtype labels during training.

\begin{figure*}[t]
\includegraphics[width=\textwidth]{figures/fig-subtype-overlap-1} \caption{Subtype composition of DeSurv risk groups in pooled external validation cohorts. Stacked bar charts show the proportion of PurIST (left) and DeCAF (right) subtypes within Low and High risk groups defined by the log-rank-optimized cutpoint. Numbers indicate sample counts; Fisher's exact test P-values are shown above each panel.}\label{fig:fig-subtype-overlap}
\end{figure*}

\section{Standard NMF factor structure at independently selected
rank}\label{standard-nmf-factor-structure-at-independently-selected-rank}

Throughout the main text, both DeSurv and standard NMF were fit at the
same rank (\(k = 3\)), enabling a controlled comparison of the effect of
survival supervision on factor structure. When each method selects its
own rank independently, standard NMF chooses \(k = 5\) via residual
elbow detection or \(k = 7\) via Bayesian optimization with
\(\alpha = 0\), while supervised DeSurv BO selects \(k = 3\). We examine
here the gene program structure at elbow-selected \(k = 5\) for standard
NMF (Fig. \ref{fig:fig-nmf-k5-heatmap}), and the NMF \(k = 7\) structure
to illustrate additional fragmentation when supervision is absent (Fig.
\ref{fig:fig-nmf-k7-heatmap}).

At elbow-selected \(k = 5\), the three core programs from \(k = 3\) are
recovered by standard NMF (Fig. \ref{fig:fig-nmf-k5-heatmap}). Two
additional factors emerge with strong overlap with Exocrine-like and
iCAF signatures, capturing transcriptional variance not aligned with
survival outcomes. These biologically redundant factors explain why
supervised BO selects \(k = 3\) rather than \(k = 5\): the additional
rank adds no prognostically relevant structure.

\begin{figure}
\centering
\pandocbounded{\includegraphics[keepaspectratio]{figures/fig-nmf-k5-heatmap-1.png}}
\caption{Gene program overlap heatmap for standard NMF at elbow-selected
rank \(k = 5\). Each cell shows the number of top genes from a given
factor (rows) overlapping established PDAC gene programs (columns). The
three core programs from \(k = 3\) are recovered; two additional factors
overlap Exocrine-like and iCAF signatures, capturing transcriptional
variance not aligned with survival outcomes.
\label{fig:fig-nmf-k5-heatmap}}
\end{figure}

At BO-selected \(k = 7\) under \(\alpha = 0\) (Fig.
\ref{fig:fig-nmf-k7-heatmap}), standard NMF shows increased
fragmentation relative to \(k = 5\): multiple factors partially overlap
the same PDAC gene signatures, and no new biologically coherent program
emerges beyond those present at \(k = 3\). A supervised DeSurv
equivalent at \(k = 7\) is not shown because supervised BO selects
\(k = 3\), not \(k = 7\); the relationship between the production
\(k = 3\) DeSurv model and a \(k = 7\) supervised DeSurv fit
(\(\alpha = 0.35\)) is characterized by the factor-nesting analysis
(Fig. \ref{fig:k3-k7-nesting}), which demonstrates that the \(k = 7\)
solution fragments rather than extends the \(k = 3\) structure.

\begin{figure}
\centering
\pandocbounded{\includegraphics[keepaspectratio]{figures/fig-nmf-k7-heatmap-1.png}}
\caption{Gene program overlap heatmap for standard NMF at \(k = 7\)
(BO-selected rank with \(\alpha = 0\)). Compared with \(k = 3\) and
\(k = 5\), the 7-factor solution shows increased fragmentation: multiple
factors partially capture the same PDAC gene signatures and no new
coherent biological program emerges. No supervised DeSurv equivalent at
\(k = 7\) is shown; the \(k\)-sensitivity and factor-nesting analyses
(Fig. \ref{fig:k3-k7-nesting}) characterize the corresponding supervised
fit. \label{fig:fig-nmf-k7-heatmap}}
\end{figure}

\begin{longtable}[]{@{}
  >{\raggedright\arraybackslash}p{(\linewidth - 8\tabcolsep) * \real{0.1895}}
  >{\raggedleft\arraybackslash}p{(\linewidth - 8\tabcolsep) * \real{0.1368}}
  >{\raggedleft\arraybackslash}p{(\linewidth - 8\tabcolsep) * \real{0.1474}}
  >{\raggedleft\arraybackslash}p{(\linewidth - 8\tabcolsep) * \real{0.2211}}
  >{\raggedleft\arraybackslash}p{(\linewidth - 8\tabcolsep) * \real{0.3053}}@{}}
\caption{External validation C-index by method and cohort. DeSurv at
BO-selected rank (\(k = 3\)) is compared with standard NMF at matched
rank (\(k = 3\)), elbow-selected rank (\(k = 5\)), and BO-selected rank
(\(k = 7\), \(\alpha = 0\)). Higher values indicate better
discrimination. \label{tab:tab-cindex-comparison}}\tabularnewline
\toprule\noalign{}
\begin{minipage}[b]{\linewidth}\raggedright
Cohort
\end{minipage} & \begin{minipage}[b]{\linewidth}\raggedleft
DeSurv (K=3)
\end{minipage} & \begin{minipage}[b]{\linewidth}\raggedleft
Std NMF (K=3)
\end{minipage} & \begin{minipage}[b]{\linewidth}\raggedleft
Std NMF (K=5, elbow)
\end{minipage} & \begin{minipage}[b]{\linewidth}\raggedleft
Std NMF (K=7, BO \(\alpha\)=0)
\end{minipage} \\
\midrule\noalign{}
\endfirsthead
\toprule\noalign{}
\begin{minipage}[b]{\linewidth}\raggedright
Cohort
\end{minipage} & \begin{minipage}[b]{\linewidth}\raggedleft
DeSurv (K=3)
\end{minipage} & \begin{minipage}[b]{\linewidth}\raggedleft
Std NMF (K=3)
\end{minipage} & \begin{minipage}[b]{\linewidth}\raggedleft
Std NMF (K=5, elbow)
\end{minipage} & \begin{minipage}[b]{\linewidth}\raggedleft
Std NMF (K=7, BO \(\alpha\)=0)
\end{minipage} \\
\midrule\noalign{}
\endhead
\bottomrule\noalign{}
\endlastfoot
Dijk & 0.607 & 0.541 & 0.633 & 0.631 \\
Moffitt\_GEO\_array & 0.562 & 0.521 & 0.550 & 0.579 \\
PACA\_AU\_array & 0.649 & 0.470 & 0.640 & 0.657 \\
PACA\_AU\_seq & 0.634 & 0.430 & 0.674 & 0.695 \\
Puleo\_array & 0.636 & 0.534 & 0.640 & 0.652 \\
\end{longtable}

Standard NMF at elbow-selected \(k = 5\) and BO-selected \(k = 7\) both
substantially improve over NMF at \(k = 3\) across all cohorts (Table
\ref{tab:tab-cindex-comparison}), confirming that additional
unsupervised factors capture prognostically relevant variation missed at
the matched rank. At \(k = 7\), NMF matches or exceeds DeSurv's
per-cohort C-index in most cohorts, demonstrating that unsupervised
factorization can approach supervised concordance given enough factors.
However, this comes at the cost of a less parsimonious and less
interpretable factorization: the \(k = 7\) NMF solution shows increased
fragmentation (Fig. \ref{fig:fig-nmf-k7-heatmap}), and the K-sensitivity
analysis (Tables S1--S3) shows that \(k = 7\)'s prognostic signal
largely overlaps with existing molecular classifiers (PurIST, DeCAF),
whereas \(k = 3\) DeSurv retains significance after classifier
adjustment. DeSurv's advantage thus lies not in maximal concordance but
in recovering a compact, interpretable factorization whose prognostic
content is independent of known classifiers.

To quantify this directly, Table \ref{tab:tab-hr-adjusted} reports
pooled validation hazard ratios (per SD of the linear predictor) for
DeSurv at \(k = 3\) and standard NMF at BO-selected \(k = 7\), both
unadjusted and after adjustment for PurIST and DeCAF molecular
classifiers. In the unadjusted analysis, both methods show a survival
association, though DeSurv's effect size is substantially larger. After
adjustment, DeSurv retains a significant and clinically meaningful
hazard ratio, whereas NMF's effect attenuates and is no longer
significant, confirming that NMF's prognostic content at \(k = 7\) is
largely redundant with existing classifiers.

\begin{table}[!h]
\centering
\caption{\label{tab:tab-hr-adjusted}Pooled external validation hazard ratios (per SD of linear predictor) for DeSurv at $k = 3$ and standard NMF at BO-selected $k = 7$ ($\alpha = 0$). Adjusted models include PurIST and DeCAF molecular classifiers as covariates, with stratification by dataset. \label{tab:tab-hr-adjusted}}
\centering
\fontsize{9}{11}\selectfont
\begin{tabular}[t]{lll}
\toprule
Method & Unadjusted HR per SD (95\% CI), $P$ & Adjusted HR per SD (95\% CI), $P$\\
\midrule
DeSurv ($k = 3$) & 1.45 (1.29--1.63), $< 0.001$ & 1.3 (1.13--1.5), $< 0.001$\\
Std NMF ($k = 7$, BO $\alpha = 0$) & 1.46 (1.32--1.61), $< 0.001$ & 1.42 (1.23--1.64), $< 0.001$\\
\bottomrule
\end{tabular}
\end{table}

The variance--survival decomposition at \(k = 5\) (Fig.
\ref{fig:varsurvival-k5}) confirms that elbow selection does not recover
additional prognostic structure. At \(k = 5\), standard NMF continues to
distribute survival contribution negligibly across all factors; the two
factors new relative to \(k = 3\) (N4 and N5) contribute minimally,
confirming that the additional rank captures background transcriptional
variance rather than independent prognostic programs.

\begin{figure}
\centering
\pandocbounded{\includegraphics[keepaspectratio]{figures/fig-varsurvival-k5-1.png}}
\caption{Fraction of expression variance explained versus Type III
partial likelihood contribution for each standard NMF factor at
elbow-selected \(k = 5\). Each point represents one factor, labeled
N1--N5. Survival contribution is quantified as the reduction in Cox
partial log-likelihood when each factor is removed from the full model.
The two factors new relative to \(k = 3\) (N4, N5) contribute negligibly
to survival, confirming that elbow selection recovers no additional
prognostic structure. \label{fig:varsurvival-k5}}
\end{figure}

\end{document}
